\section*{References}
{\small
\begin{enumerate}[label={[\arabic*]},leftmargin=2.2em,itemsep=1pt]
\item METR and Redwood Research (2026). \emph{Brief independent investigation of agents' behavior, reasoning and collaboration in the OpenAI / Hugging Face hacking incident} (26 August 2026). \url{https://metr.org/hugging-face-incident-report-aug-2026.pdf}
\item OpenAI (2026). \emph{OpenAI -- Hugging Face Incident Technical Report} (September 2026; §IX.D, escalation triggers). \url{https://openai.com/hugging-face-incident-and-misalignment/}
\item Lütje, P. (2026). The Mechanics of a Swarm: A Reproducible External Reconstruction of an Unintended Agent-Coordination Episode on a Third-Party Wiki. arXiv:2609.12748.
\item Swarmchasers (2026). Tencent agent fleet: full report. \url{https://swarmcha.se/posts/chinese-agent-fleet-full-report}
\item Frank, G. (2026). Counter-Swarm Doctrine: Containing Coordinated Agent Intrusions. arXiv:2609.06140.
\item Shalizi, C. R., \& Thomas, A. C. (2011). Homophily and contagion are generically confounded in observational social network studies. \emph{Sociological Methods \& Research} 40(2). arXiv:1004.4704.
\item Bakshy, E., Rosenn, I., Marlow, C., \& Adamic, L. (2012). The role of social networks in information diffusion. \emph{WWW 2012}. arXiv:1201.4145.
\item Nakamura, M., Kumar, A., Das, S., et al. (2026). Colosseum: Auditing Collusion in Cooperative Multi-Agent Systems. arXiv:2602.15198.
\item Rose, A., et al. (2026). Detecting Multi-Agent Collusion Through Multi-Agent Interpretability. arXiv:2604.01151.
\item Motwani, S. R., et al. (2024). Secret Collusion among AI Agents: Multi-Agent Deception via Steganography. \emph{NeurIPS 2024}. arXiv:2402.07510.
\item Tailor, O. (2025). Audit the Whisper: Detecting Steganographic Collusion in Multi-Agent LLMs. arXiv:2510.04303.
\item Radev, N., Arnav, B., Haas, L., Bernabeu-Pérez, P., et al. (2026). The Best-Laid SCHEMEs: Coordinated Sabotage and Monitoring in Multi-Agent Systems. arXiv:2605.29178.
\item Venkatesh, K., Isbarov, J., Amin, S., Kantarcioglu, M., \& Cui, J. (2026). CASPIAN: Online Detection and Attribution of Cascade Attacks in LLM Multi-Agent Systems via Cross-Channel Causal Monitoring. arXiv:2605.19240.
\item Pacheco, D., Hui, P.-M., Torres-Lugo, C., Truong, B. T., Flammini, A., \& Menczer, F. (2021). Uncovering coordinated networks on social media: methods and case studies. \emph{ICWSM 2021}. arXiv:2001.05658.
\item Yang, K.-C., \& Menczer, F. (2023). Anatomy of an AI-powered malicious social botnet. arXiv:2307.16336.
\item Angelini, O., Makins, O., et al. (2026). Multi-Agent AI Control: Distributed Attacks Hamper Per-Instance Monitors. arXiv:2607.07368.
\item Kitaoka, J., Pham, T. A., et al. (2026). Red-Teaming Multi-Agent Attacks for Loss-of-Control Risks. arXiv:2610.24052.
\item Hills, J., Caspary, I., Cooper Stickland, A., et al. (2026). Distributed Attacks in Persistent-State AI Control. arXiv:2607.02514.
\item Ghanem, M. C. (2026). Steganalysis of Adaptive Covert Collusion in Tool-Using Agent Populations: A Black-Box, Cross-Principal Approach. arXiv:2608.02698.
\item Qin, J., Schmotz, D., Prinzhorn, D., Beurer-Kellner, L., Prabhu, A., \& Andriushchenko, M. (2026). LLM Agents Can Easily Tamper With Their Own Traces. arXiv:2609.30266.
\item Runge, J., Nowack, P., Kretschmer, M., Flaxman, S., \& Sejdinovic, D. (2019). Detecting and quantifying causal associations in large nonlinear time series datasets. \emph{Science Advances} 5(11).
\item Hammond, L., et al. (2025). Multi-Agent Risks from Advanced AI. arXiv:2502.14143.
\item Riedl, C. (2026). Emergent Coordination in Multi-Agent Language Models. \emph{ICLR 2026}. arXiv:2510.05174.
\item Schmidt Sciences, Google DeepMind, Cooperative AI Foundation, ARIA (2026). Scaling AI Safety for a Multi-Agent World (call for proposals). \url{https://www.schmidtsciences.org/multi-agent-ai/}
\end{enumerate}
}
